\scanpage{17}
voisinage.
\figure{\begin{tikzpicture}[x=1cm,y=1cm]\draw[gray](0,.37)--(.7,.5);\draw[gray](.05,.3)--(.68,.4);\draw[gray](.25,0)--(.27,.43);\draw[gray](.61,.06)--(.64,.46);\end{tikzpicture}}

Nous désignerons par
\[
F=F_\varphi=\underline{\operatorname{Hom}}_S(X,Y;\varphi)
\]
le sous-foncteur précédent de $\underline{\operatorname{Hom}}_S(X,Y)$, dont la valeur en $S'$ est l'ens des \unsure{$S'$-homs} $u:X'-Y'$ tels que, posant $v=u\varphi'$, et considérant $M_v^{\unsure{\cdot|\prime}},M_u,M_v\longrightarrow M_u$ correspondant, $M_v\longrightarrow M_u$ soit surjectif.

\uline{Théorème} $\underline{\operatorname{Hom}}_S(X,Y;\varphi)\longrightarrow\underline{\operatorname{Hom}}_S(Z,Y)$ est représentable par morphismes non ramifiés.
\figure{\begin{tikzpicture}[x=1cm,y=1cm]\draw[gray](0,0)..controls(.4,.45)and(.7,.25)..(.95,.04);\draw[gray](2.3,.03)..controls(2.4,.45)and(3,.32)..(3.45,.08);\end{tikzpicture}}

\uline{Corollaire} Si $\underline{\operatorname{Hom}}_S(Z,Y)$ est représentable, alors $\underline{\operatorname{Hom}}_S(X,Y;\varphi)$ l'est également, et le morphisme $\underline{\operatorname{Hom}}_S(X,Y;\varphi)-\underline{\operatorname{Hom}}_S(Z,Y)$ est non ramifié
\figure{\begin{tikzpicture}[x=1cm,y=1cm]\draw (0,0)--(-.04,1.2);\end{tikzpicture}}

\uline{Principe de la démonstration} On peut supposer donné un $v:\operatorname{Hom}_S(Z,Y)$, et on doit représenter le sous-foncteur de $\operatorname{Hom}_S(X,Y)$ dont la valeur en $S'$ est l'ens des $S'$-morphismes $u':X'-Y'$ tels que 1°) $u'\circ\varphi'=v'$ 2°) La condition $M_{v'}\longrightarrow M_{u'}$ surjectif. On applique le critère général de représentabilité \inserted{pour} des \deleted{foncteurs}
\margin{N.B. Réduction au cas $S$ noeth. immédiate \unsure{sans doute sur} foncteurs rel. représentables}
\mspage{16}
